% 第 21 章　电流和直流电路
\chapter{电流和直流电路}

\step{反常识开场}

晚上回家，你在门口按下开关，天花板上的灯几乎在同一瞬间亮了。这件事平凡到没有人会多看一眼。但本章后面我们会做一道简单的算术题，结论会让这个瞬间变得可疑：在普通铜导线里，自由电子沿着导线“前进”的平均速度大约只有每分钟几毫米。也就是说，开关附近的那个电子，就算灯泡亮一整夜，它也走不到灯泡那里去。既然“搬运电的家伙”动得这么慢，灯凭什么立刻就亮？

先别急着回答。再看第二个场景。高压输电线上的电压可以高达几十万伏，远远超出家庭插座的二百二十伏。可是鸟常常一只挨一只地站在裸露的高压线上，悠闲地梳理羽毛，什么事也没有。反过来，家里插座只有二百二十伏，说明书和大人却反复警告：绝对不能碰。如果危险程度由电压大小决定，鸟早就该出事了。

第三个场景藏在你的口袋里。给手机充电时摸摸充电线，有时会感到微微发热；同样是充电，粗一些、短一些的原装线往往比又细又长的杂牌线充得快。线里面到底发生了什么，让“粗细”和“长短”都成了变量？

灯、鸟、充电线，三件事指向同一个问题：电流究竟是什么东西在动，它以多快的速度动，又是什么在推动它动？这一章我们就把这条回路拆开看清楚。

\step{先猜一猜}

在往下读之前，请把你的直觉押在桌面上。对下面几个问题，先写下你的答案，不许翻页偷看。

第一问：一节电池接一个小灯泡，电流从电池正极出发，经过灯泡，回到负极。你猜：在“灯丝”这个地方和“回到电池的导线”这个地方，两处每秒钟通过的电荷一样多，还是不一样多？很多人隐隐觉得，灯泡“消耗”了电，所以回来的那一路应该弱一些。

第二问：把两个完全一样的小灯泡一前一后接在同一节电池上（串成一条线），再和“只接一个灯泡”比，你猜每个灯泡会更亮、更暗，还是不变？如果把它们并排各接各的（并联）呢？

第三问：鸟站在高压线上没事。你猜原因是“鸟的脚不导电”“鸟离地面远”，还是别的什么？

第四问：一根导线越细、越长，电流通过它是更容易还是更难？你大概觉得更难——那“难多少”？和长度成正比，还是和长度的平方成正比？

这些猜测不需要任何公式，只需要你对世界的朴素感觉。物理学家的感觉也是这样开始的；区别只在于，他们接下来会设计实验，逼世界表态。

\step{动手实验}

\begin{experiment}[铅笔芯变阻器]
你需要：一节五号电池（或两节串联）、一个手电筒用的小灯泡（或一颗发光二极管加一百欧左右的限流电阻）、几根导线（可用去掉绝缘皮的回形针代替）、一支木头铅笔、一把小刀。

第一步：把小刀削开铅笔的一侧，露出整条石墨笔芯。石墨是能导电的，这是本实验的关键道具。

第二步：用导线把电池、灯泡连成回路，中间留一个缺口。缺口的两个线头，一个夹在笔芯的一端，另一个线头做成可以在笔芯上滑动的夹子。

第三步：让滑动夹从笔芯的一端慢慢滑向另一端。你会看到灯泡从暗到亮连续变化——接入回路的笔芯越短，灯越亮；笔芯越长，灯越暗。你已经亲手做出了一个滑动变阻器，收音机的音量旋钮、台灯的无级调光，原理和它一模一样。

第四步（可选，注意安全）：把滑动夹移到灯最亮的位置，保持通电一分钟左右，再用手指轻触笔芯。它会微微发热。电能没有消失，它在笔芯里变成了热。
\end{experiment}

还有一个不用任何器材的“人体实验”，留到人多的时候做：请十几个同学在教室里手拉手围成一个大圈，每个人扮演一个“电荷”。让“电池”同学匀速地推身边的人，推力沿圆圈依次传下去。注意两件事：其一，每个人只挪了一小步，但“有人在动”这个消息几乎瞬间传遍了整圈；其二，圈上每个位置，“单位时间经过的人”必然一样多——否则人就要在某个地方堆积起来了。记住这两个观察，它们后面会变成两条定律。

最后补一个纯观察任务：找出你身边的三样电器——手机充电器、电热水壶、台灯——读一读它们的铭牌，记下“伏”和“瓦”两个数。用本章稍后的公式心算一下每样东西的工作电流，再回答：哪一样最需要粗导线？哪一样如果内部短路后果最严重？把答案记下来，读完本章再回来核对。

\step{旧模型遇到麻烦}

大多数人心里有一个默认模型：“电”是电池里装着的一种东西，像水箱里的水；打开开关，它顺着导线流出来，流过灯泡时被“用掉”一部分，剩下的流回电池；电池用光了，就是“电”流完了。

这个模型很顺嘴，但它经不起三记敲打。

第一记敲打来自实验事实。如果你把电流表（或者灵敏的检测手段）串在回路的不同位置——电池旁边、灯泡前面、灯泡后面、回到电池的最后一截导线——会发现一个朴素却惊人的结果：在只有一条通路的回路里，处处每秒钟通过的电荷量相同，分毫不差。灯丝前面是多少，灯丝后面就是多少。如果“电被灯泡用掉了”，回来那一路理应减少；可测量说没有。灯泡并没有吞掉电荷，它从电荷身上拿走了另一样东西。

第二记敲打来自开场的算术（下文我们会真的算出来）：电子在导线里的平均前进速度只有每小时零点几米量级。如果灯泡亮是靠“开关处的电跑过来说一声”，那灯要等上几个小时。可灯是立刻亮的。“某个东西跑得飞快”和“载流粒子本身在爬行”，这两者必须同时成立，旧模型装不下这个矛盾。

第三记敲打来自鸟。高压线上的鸟脚下的电压是几十万伏，可它毫发无损；而一节一点五伏的电池，用一根铜丝直接连通正负两极，铜丝却能在几秒内烫到让你缩手，电池也很快报废。低电压在这里反而制造了危险。这说明“电压高”本身既不是危险的充分条件，也不是必要条件。决定后果的，是“有多少电荷、以多大的能量代价、穿过了身体的哪一部分”。旧模型里“电压越高越吓人”的一句话，解释不了这些细节。

顺带说一句常被误用的话：“电被电池用完了”。电（电荷）没有被用完——回路里的电荷自始至终在循环，一进一出严格相等；被用完的是电池里储存的化学能。真正被消耗、被转移的是能量，不是电荷。这一条我们会在本章反复使用。

旧模型该退休了。但我们不把它完全扔掉——“像水流”这个感觉在很多时候依然好用。我们只是要弄清楚：这个比喻像什么，又不像什么。

\step{发明新概念}

历史上，物理学家解决这些麻烦的办法，是发明一组精确的新词。我们按需要的顺序把它们造出来。

第一个新概念：电流。回到人体圆圈的实验。描述“流量”最自然的办法，不是问“某个人跑多快”，而是问“每秒钟有多少人通过某个位置”。对电荷也一样：任取导线的一个横截面，数一数每秒钟有多少电荷穿过它。于是我们定义

\[
I=\frac{\Delta Q}{\Delta t},
\]

即在时间 \(\Delta t\) 内穿过某截面的总电荷 \(\Delta Q\) 除以这段时间。电流 \(I\) 是电荷的流率，单位是安培，一安培就是每秒一库仑。注意这个定义里根本没有出现“电子的速度”：同样的电流，可以由很多慢速的电荷组成，也可以由少量快速的电荷组成，就像同样大的客流量，可以由人山人海缓步挪动造成，也可以由稀疏人群快速通过造成。这就是“电流不是电子运动速度”的准确含义。

第二个新概念：漂移速度。导线里的自由电子其实一直在做剧烈的无规则热运动，速率高达每秒十万米的量级，但方向乱七八糟，平均下来哪儿也没去——这就是没有电流时的状态：喧嚣，但没有净流动。接通电路后，电子在热运动的混乱之上，叠加了一个极其缓慢的定向“蠕行”，平均每分钟只前进几毫米。这个平均的定向速度叫漂移速度。电流、信号、热运动，是三件不同的事：热运动是每只脚的乱走，漂移是整个队伍的平均挪移，而“开关闭合了”这个消息的传播，又是另一回事。

第三个新概念：电场推动，消息以近光速传播。开关一合上，电池两端的电势差立刻在整条回路里建立起电场。这个建立过程不是由电子从开关“跑”过去完成的，而是电荷分布以接近光速的速度在导线中重新调整——人浪实验里，你不需要等第一个人走到终点，推力沿着队伍依次传递，整圈人几乎同时开走。电场像什么？像一根已经灌满水的长水管：阀门一开，出口立刻出水，因为水早就充满了管子，压力的变化以水中的声速传遍整根管子。它不像什么？不像空水管等水从源头流过来。导线里自由电子本来就无处不在，灯泡灯丝里也有，所以电场一到，灯丝里的电子当场就开始定向移动，灯立刻亮。开场的第一个悬念，到这里已经能回答了。

第四个新概念：电压是“每份电荷的能量台阶”。第 20 章已经把电势差讲清楚了，这里只强调它在电路里的角色：电压不是电流的“压力”，而是每单位电荷跨过两点之间时会获得或付出的能量。一盏灯两端的电压是一伏，意味着每一库仑电荷穿过灯丝，就要交出一焦耳的能量，变成光和热。请注意三个词的区别：电势是某个点相对于约定零点的“高度”；电势能是某个电荷处在那个高度时拥有的能量，等于电荷量乘电势；电压是两点之间高度的差。本章用到的永远是第三个——回路里真正干活的，从来是“差”，不是某一点的绝对值。鸟的秘密也藏在这里，稍后揭晓。

第五个新概念：电阻。同样的电池，接粗短的铜线，电流很大；接细长的石墨笔芯，电流小得多。材料、粗细、长短共同决定了一个量：在推动电荷这件事上，这段导体有多“为难”。我们把电压和电流的比值定义为这段导体的电阻

\[
R=\frac{V}{I},
\]

单位是欧姆。电阻大，不是“把电流挡住了”，而是“同样的电压只能推动更小的电流”——或者说，要推动同样的电流，必须付出更大的电压。铅笔芯实验里滑动的夹子，改变的就是接入回路的电阻。

第六个新概念：电池是一台化学“泵”。电荷在回路里循环，每绕一圈，在灯丝和导线里把能量交出去，总得有个地方把能量补回来，否则流动很快就会停。这个角色由电池扮演：电池内部的化学反应持续地把正电荷从低电势一端“搬”到高电势一端，像水泵把水从低处抽到高处，维持住两极之间的电势差。化学反应每搬运一库仑电荷所提供的能量，就是电池的电动势——名字里虽然有个“力”字，它的单位却是伏特，它实际上是能量，不是力。电池没电了，不是电荷用完了，而是能驱动这场搬运的化学物质消耗完了。

\step{一句话模型}

\begin{oneline}
电路是一条首尾相接的传送带：电池用化学能把电荷不断泵上“能量高坡”，电荷沿闭合回路流动，把能量沿途交给灯泡、电阻和导线，再回到电池重新上坡；回路里任何截面单位时间通过的电荷都相等（电荷不会凭空增减），绕回路一整圈能量的升降总和为零（能量不会凭空增减）——前一条叫基尔霍夫电流定律，后一条叫基尔霍夫电压定律。
\end{oneline}

\step{数学翻译器}

现在把上面的直觉翻译成可以计算的式子。每出现一个式子，我们就立刻用极端情形“拷问”它一次。

\textbf{第一式：电流与漂移速度。}设导线横截面积为 \(A\)，单位体积里有 \(n\) 个自由电荷，每个电荷为 \(q\)，以漂移速度 \(v\) 前进。在时间 \(\Delta t\) 内，能穿过某截面的电荷恰好是体积 \(Av\Delta t\) 里的全部自由电荷，即 \(\Delta Q=nqAv\Delta t\)，于是

\[
I=nqAv.
\]

来看真实数量级。家庭电路常用直径约一毫米的铜线，截面积 \(A\) 约为 \(8\times10^{-7}\) 平方米；铜的自由电子数密度 \(n\) 约为 \(8.5\times10^{28}\) 每立方米；电子电荷 \(q\) 约为 \(1.6\times10^{-19}\) 库仑。取一个不小的电流 \(I=1\) 安培代入，解出

\[
v=\frac{I}{nqA}\approx\frac{1}{8.5\times10^{28}\times1.6\times10^{-19}\times8\times10^{-7}}\approx9\times10^{-5}\ \text{米每秒}.
\]

也就是每秒零点零九毫米，一分钟约五毫米。检查特例：电流为零时漂移速度为零（热运动还在，只是没有净流动），合理；截面积越大，同样电流下漂移越慢，合理——粗水管里水流更从容。这就是开场算术的来源：灯立刻亮，靠的不是电子赶路，而是电场以接近光速建立。

顺便拆穿一个商场里的单位游戏。充电宝爱标“一万毫安时”。毫安时是电流乘时间，按 \(I=\Delta Q/\Delta t\) 反推，它是电荷的单位：一万毫安时等于三万六千库仑，约合 \(2\times10^{23}\) 个电子的电荷量。但电荷量还不是能量——同样多的电荷，在三点七伏的电芯里和在五伏的输出端，携带的能量不同。所以比较充电宝，要看“瓦时”（能量），而不是只看“毫安时”（电荷）。这正是本章反复强调的：回路里守恒地循环的是电荷，真正被消耗、被计价的是能量。

\textbf{第二式：欧姆定律。}对许多导体，在温度变化不大时，测出的电流与所加电压成正比：

\[
V=IR.
\]

它回答“要推动电流 \(I\) 通过电阻 \(R\)，需要多大的电压”。立刻拷问它：若 \(R\) 趋于零（用铜丝直连电池两极，即短路），公式预言电流趋于无穷大——现实中当然达不到无穷，因为电池内部也有电阻、铜丝也有电阻，但电流足以让铜丝几秒内发烫变红，这正是为什么绝不能短接电池，也是保险丝存在的理由。若回路断开，相当于 \(R\) 无穷大，预言电流为零——开路无电流，合理。注意，欧姆定律不是一个像牛顿定律那样的普适定律，它是对一大类材料在常用条件下的经验总结，电阻 \(R\) 本身可能随温度、电流变化；它的边界我们下一节再谈。

\textbf{第三式：串联与并联。}先把语言统一。工程师不画实物，画电路图：电池用一长一短两条平行线表示（长线是正极），开关用一个可以合上的缺口表示，灯泡和电阻用方框或圆圈表示，导线拉成直线。下面这张图就是本章反复讨论的“一节电池、一个开关、一盏灯”：

\begin{center}
\begin{tikzpicture}[>=Stealth,scale=1.0]
  % 导线回路
  \draw (0,0) -- (6,0) -- (6,1.15);
  \draw (6,1.85) -- (6,3) -- (3.35,3);
  \draw (2.65,3) -- (0,3) -- (0,1.7);
  \draw (0,1.3) -- (0,0);
  % 电池（左侧）：长细线为正极，短粗线为负极
  \draw[thick] (-0.5,1.7) -- (0.5,1.7);
  \draw[line width=2.5pt] (-0.25,1.3) -- (0.25,1.3);
  \node[left] at (-0.5,1.7) {$+$};
  \node[left] at (-0.5,1.3) {$-$};
  \node[left] at (0,0.7) {电池};
  % 开关（上方）
  \fill (2.65,3) circle (1.5pt);
  \fill (3.35,3) circle (1.5pt);
  \draw (2.65,3) -- (3.3,3.45);
  \node[above] at (3,3.45) {开关};
  % 灯泡（右侧）：圆圈加叉
  \draw (6,1.5) circle (0.35);
  \draw (5.75,1.25) -- (6.25,1.75);
  \draw (5.75,1.75) -- (6.25,1.25);
  \node[right] at (6.45,1.5) {灯泡};
  % 电流方向箭头
  \draw[->,thick] (1.2,0) -- (2.2,0);
  \node[below] at (1.7,0) {$I$};
\end{tikzpicture}
\end{center}

约定电流 \(I\) 的箭头从电池正极出发，沿外电路回到负极（这是历史约定，金属里实际漂移的电子方向相反，效果等价）。有了这张图，“串联”就是所有元件串在唯一的回路上，“并联”就是几个元件各自架在同一对节点之间。

两个电阻首尾相接（串联），流过它们的电流必为同一个 \(I\)（电荷守恒，没别处可去），总电压是两段电压之和，所以

\[
R_{\text{串}}=R_1+R_2.
\]

每个电阻分到的电压与其阻值成正比：\(V_1=V\,R_1/(R_1+R_2)\)，这就是分压。两个电阻两端接在同一对节点上（并联），两端电压必为同一个 \(V\)，总电流是两路之和，所以

\[
\frac{1}{R_{\text{并}}}=\frac{1}{R_1}+\frac{1}{R_2}.
\]

每路分到的电流与其阻值成反比，这就是分流。拷问它们：串联里若有一个电阻为零，总电阻不退化、公式仍成立；若有一个断开（视为无穷大），整条路电流为零——老式节日灯串几十个小灯泡串联，一个烧断整串熄灭，就是这个道理。并联里若某一支断开，其余支路照常工作——家里电灯、冰箱、插座全是并联在市电上的，关掉台灯不影响冰箱，而且每加一个用电器，从电源看到的总电阻反而变小、总电流变大。这就是“先猜一猜”第二问的答案：两个同样的灯泡串联，总电阻加倍，电流减半，每个灯都明显变暗；并联则每个灯的电压不变，亮度基本不变（理想电池下），但电池的负担加倍。

\textbf{第四式：基尔霍夫两条定律。}串联并联公式其实只是两条更普遍定律的特例。

电流定律：在电路的任意一个节点，流进的电流之和等于流出的电流之和。它不是新定律，就是“电荷不能在节点处堆积或凭空产生”——这正是人体圆圈实验里“每个位置流量相同”的一般化。

电压定律：沿任何一条闭合回路绕一圈，所有电压的升降之和为零。它也不是新定律，就是“每单位电荷绕一圈回到出发点，电势能必须恢复原值”，背后是能量守恒。

有了这两条，任何只用电池和电阻搭成的直流电路，原则上都能算：给每条支路设一个未知电流，列方程，解方程。

\begin{mathbridge}
来看一个串联并联公式处理不了的情形：两节电池 \(V_1=6\) 伏、\(V_2=3\) 伏，和三个电阻 \(R_1=2\) 欧、\(R_2=4\) 欧、\(R_3=2\) 欧接成一个双回路网。设两条支路电流 \(I_1\)、\(I_2\)，第三条支路由电流定律自动给出 \(I_1+I_2\)（或差，取决于方向约定）。对两个回路分别用电压定律：

\[
6=2I_1+2(I_1+I_2),\qquad 3=4I_2-2(I_1+I_2).
\]

这是两个二元一次方程，解出 \(I_1\)、\(I_2\) 即可。如果某个解出来是负数，只说明实际方向与你假设的箭头相反——方向反转的情形，公式用负号自动处理，这也是对模型的一种“拷问”。更重要的信息是：基尔霍夫定律把“解电路”变成了解线性方程组。电路再大，也只是方程更多，交给计算机即可——电路仿真软件每天做的就是这件事。
\end{mathbridge}

\textbf{第五式：电功率与焦耳热。}每库仑电荷跨过电压 \(V\) 交出能量 \(V\) 焦耳，每秒通过 \(I\) 库仑，所以用电器每秒获得的能量，即功率

\[
P=VI.
\]

若用电器服从欧姆定律，还可以写成 \(P=I^{2}R\) 或 \(P=V^{2}/R\)。这些能量在电阻性元件里几乎全部变成热，这就是焦耳热。来一道真实数据的估算：一个标着“220 伏、2000 瓦”的电热水壶，工作电流 \(I=P/V\approx9.1\) 安培，电阻 \(R=V^{2}/P\approx24\) 欧。把一点五升水从二十摄氏度烧开大约需要五十万焦耳，两千瓦的功率意味着约二百五十秒，四分多钟——和你家厨房的实际体验吻合。耗电约零点零八度，按每度电几角钱算，烧开一壶水的电费只有几分钱。另一个估算：手机充电器输出五伏两安，功率十瓦；如果你用又细又长的劣质线，线的电阻可能达到一欧，单根导线上的发热功率就有 \(I^{2}R\approx4\) 瓦量级，同时分走的电压让手机端电压不足，充电变慢、线身发热——开场第三个悬念解决。

\textbf{一个让直觉摔跤的反例。}看 \(P=I^{2}R\) 和 \(P=V^{2}/R\)：前一个说电阻越大发热越多，后一个说电阻越大发热越少。到底谁对？这正是直觉容易翻车的地方。答案是：都对，但前提不同。电流被强制保持不变的场合（比如串联回路里），谁的电阻大谁发热多——所以接头松动处电阻变大，反而最容易烧坏；电压被强制保持不变的场合（比如并联在市电上的用电器），谁的电阻大谁功率小——所以“大功率电器”恰恰电阻小。一句话：比较之前，先问什么是被固定住的。这也是远距离输电要用高压的原因：输送功率 \(P=VI\) 一定时，电压越高电流越小，而线路发热 \(I^{2}R\) 按电流的平方下降——用几十甚至上百万伏输电，不是为了吓唬人，是为了少发热。

另一个有真实数量级的例子是汽车电瓶：区区十二伏，启动发动机时却能输出几百安培，功率达到几千瓦——好在只维持几秒钟。电瓶的“本事”不在电压高，而在内阻极小，放得出大电流。这也再次说明：离开电流谈电压，或者离开电压谈电流，都描述不了能量的吞吐。

焦耳热是祸害也是工具，差别只在于你把它放在哪里。电饭锅、电热水壶、电暖器，就是让电流老老实实通过一根发热丝；而家庭电路的安全装置，全是在防止焦耳热出现在错误的地方。保险丝是串在回路里的一小段低熔点金属：电流过大时，它先于任何导线发热熔断，用自我牺牲切断回路；空气开关做同样的事，但用电磁或热胀机构“跳闸”，可以复位。地线是另一条思路：洗衣机、冰箱的金属外壳接一根通往大地的导线，一旦内部漏电使外壳带电，电流会经电阻很小的地线流走，而不是等你的手来提供通路。漏电保护器更聪明：它时刻比较火线出去和零线回来的电流是否相等——正常情况下二者必须相等（基尔霍夫电流定律！），一旦差值超过几十毫安，说明有电流“抄近路”漏掉了，可能正穿过某个人的身体，它便在零点一秒内断电。你看，连保命的设备，工作原理都是本章那两条守恒律。人体本身的危险阈值也是真实数据：几毫安只是麻，几十毫安穿过心脏区域就可能致命——所以安全电压的说法才以干燥环境下三十六伏以下为参考。

\step{模型的边界}

每个模型都有它的活动范围。这一章的模型有四条主要边界。

第一，欧姆定律是经验规律，不是宇宙法则。金属在温度变化不大时服从它；但灯丝从室温烧到两千多摄氏度，电阻会涨好几倍——所以灯泡在刚通电的一瞬间电流最大，烧断往往发生在开灯那一刻。把电流和电压的关系画成图，欧姆型导体是一条过原点的直线，灯丝则是一条随电压升高而变弯的曲线：

\begin{center}
\begin{tikzpicture}[scale=1.6,>=Stealth]
  \draw[->] (0,0) -- (2.6,0) node[below] {$U$};
  \draw[->] (0,0) -- (0,1.8) node[left] {$I$};
  \draw[thick,blue!60!black] (0,0) -- (2.2,1.35) node[right,align=left] {欧姆型导体\\（直线）};
  \draw[thick,red!60!black] (0,0) .. controls (0.7,0.65) and (1.5,0.8) .. (2.2,0.9) node[right,align=left] {灯丝\\（变弯）};
\end{tikzpicture}
\end{center}

曲线变弯的含义是：电压越高，再增加同样的电压，电流的增加变少了——电阻在悄悄变大。二极管更极端：正向几乎畅通，反向几乎断绝，它的电压电流关系根本不是一条直线。石墨、电解质溶液、人体，各有各的脾气。用 \(V=IR\) 之前，先问对象是不是“欧姆型”的。

第二，本章的“电压”概念依赖一个不声不响的假设：电场是保守场，绕回路一圈电势差严格为零。对恒定电流和缓慢变化的电路，这个假设极好。但第 23 章你会看到，变化的磁场会产生绕圈不回到零的电场——把一个回路套在变化的磁场外面，“两点之间电压”甚至可以依赖你测量时导线绕行的路径。到那时，“电压”这个词要升级成“电动势”和“环路积分”才说得清楚。本章的电压定律，是那个更普遍规律在恒定条件下的样子。

第三，我们把导线当成没有电阻的理想连接线。普通实验里这很好；但当电流很大或导线很长时——比如从电厂到你家的输电线——导线电阻必须算进去，充电线发热的例子已经说明了这一点。

第四，漂移速度的图像是统计平均。每个电子的真实轨迹是一团乱麻上叠着缓慢蠕行，不是整齐的行军。另外，“电子带着负电向左”和“正电荷向右”在电流效果上等价，历史约定把电流方向定为正电荷运动方向，与金属中电子的实际漂移方向相反——这只是约定，不是错误，但读旧书和新书时要记得它。

第五，我们把电池当成电压恒定的理想泵。真实电池内部也有电阻（内阻）：输出电流越大，内阻分走的电压越多，对外提供的端电压越低。新电池内阻小，亏电的旧电池内阻大——这就是为什么旧电池“测电压还行，一带负载就趴下”：空载时电压表的读数接近电动势，一接上用电器，内阻立刻把端电压拉下来。手机在寒冷天气里掉电快、突然关机，部分原因也是低温让电池内阻变大。

\begin{mathbridge}
为什么金属会服从欧姆定律？一个粗粒度的微观图像是：电子在电场下被加速，但跑不了多远就撞上晶格的振动和缺陷，把定向运动的动能交出去，然后重新加速。设两次碰撞之间的平均自由时间为 \(\tau\)，电子在每次“加速—撞停”循环里获得的平均定向速度正比于电场 \(E\)，于是漂移速度 \(v\propto E\)，电流密度 \(J=nqv\propto E\)，写成

\[
J=\sigma E,
\]

比例系数 \(\sigma\) 叫电导率。欧姆定律 \(V=IR\) 就是这个微观关系对整根导线的宏观表达：电阻 \(R=L/(\sigma A)\)，正比于长度、反比于截面积——回答“先猜一猜”第四问：同样材料，电阻与长度成正比，不是平方。温度升高，晶格振动加剧，碰撞更频繁，\(\tau\) 变短，电阻变大——灯丝的脾气也有了出处。这个图像不严格（真正严格的理论要请出量子力学），但方向正确。
\end{mathbridge}

\begin{challenge}
再往下追问一层：电流在导线里流动时，电场在哪里？初学者的图像是“电场在导线内部推着电子走”，这只说对了一半。电磁场的完整理论（坡印廷矢量，第 25 章会正式讨论电磁波的能量流）告诉我们：能量并不是“在导线里面跟着电子一起流”的，而是从电池出发，通过导线周围空间的电磁场传输到用电器，导线更像引导能流的轨道。一个能检验直觉的推论：用一对平行导线给远处的灯供电，能流主要发生在两根导线之间的空间里，而不是铜的内部。铜里的电子以每分钟几毫米蠕行，而能量以接近光速从场外送达——开场“灯立刻亮”的完整回答，要到场的语言里才算说完。本章的回路定律仍然好用，因为在恒定情形下，场的描述和回路的描述严格一致；但你现在知道了，回路语言是“场语言”在低频世界的速记。
\end{challenge}

\step{回头解释开场现象}

现在清点三个开场悬念。

灯为什么立刻亮？因为灯丝里本来就有自由电子，开关闭合后，电场以接近光速在回路中建立，灯丝处的电子当场开始漂移，电流当场通过灯丝。没有任何粒子需要从开关“跑到”灯泡。这就像已经灌满水的水管，阀门一开，出口立刻出水——它不像空管子等水从水龙头爬过来。记住这个比喻的边界：水管类比能解释“立刻有水”和“处处流量相等”，但它解释不了为什么“压力”恰好按电阻比例分配，也装不下电场在空间中的传递——比喻到此为止，公式接着走。

鸟为什么没事？因为伤害来自“穿过身体的电流”，而电流由电势差驱动。鸟的两只脚站在同一根导线上相距几厘米，导线的电阻极小，这两只脚之间的电势差只有万分之一伏量级，穿过鸟身体的电流微乎其微。高压的“高”是相对于地面而言的；鸟悬浮在电线上，全身几乎处于同一个电势，没有“差”，就没有驱动。人也一样：如果只碰一根线且与大地、与其他导线绝缘，理论上不会触电——但现实中人几乎总是通过地板、墙壁、潮湿的皮肤与地面相连，二百二十伏的市电一端就接在地上，于是身体成了通路，几十毫安的电流穿过心脏区域就足以致命。危险的不是“电压高”这三个字，而是“你的身体恰好接在两个有电势差的地方”。

充电线为什么发热、为什么细长线充电慢？因为线有电阻，\(I^{2}R\) 变成热，\(IR\) 分走了本该给手机的电压。同样的电流，线越细越长，电阻越大，发热越多、压降越大。厂商把快充线做粗，就是在压低这个电阻。

三个悬念，一个模型，全部收回。这就是检验：模型不只是“能算出数”，还要能把最初让你困惑的现象一件一件认领回去。

\step{脑洞实验与挑战题}

\begin{thought}[从地球到月球的电路]
假想拉一对导线，从地球连到三十八万千米外的月球，在月球上接一盏灯，在地球上装一个开关和电池。合上开关，灯多久亮？直觉可能说“立刻”，也可能说“要等电子爬过去”——两个都错。电场建立的消息以接近光速传播，约每秒三十万千米，所以灯大约在一秒多以后亮：不慢到几小时，也快不过光速。再想一步：如果用一根从地球拉到月球的“棍子”捅一下开关，或者对着月球喊一声让声控开关闭合呢？棍子上的推动以材料中的声速传播（每秒几千米），声音根本传不过真空。三种“发信号”的方式，速度天差地别，但没有任何一种能快过光。这个思想实验把本章的“信号速度与载流子速度之别”推到极致，也预先碰到了第 28 章将要正式面对的规则：任何影响、任何消息，都不能超光速传播。
\end{thought}

\begin{puzzles}
\item 一节理想电池（无内阻）先接一个小灯泡，再并联接上第二个一模一样的灯泡。第二个灯泡接通的瞬间，第一个灯泡的亮度变不变？电池的负担变不变？换成一节有明显内阻的真实电池，结论会怎样修改？提示：先问“灯泡两端的电压由什么决定”；内阻相当于在回路里串了一个固定电阻，电流越大它分走的电压越多。
\item 老式灯串把几十个小灯泡串联，一个烧断整串全灭。有人想了个办法：在每个灯泡两端并联一根“平时绝缘、灯丝烧断后被击穿导通”的细丝。从此一个灯坏了，其余的还亮。请用本章模型说明这个巧妙设计的代价——剩下的灯泡会怎样变化？提示：坏灯位置从“断开”变成“近似短路”，总电阻减小，总电流增大，其余每个灯分到的电压略微上升、变得更亮也更危险；坏的越多，剩下的越亮，这是恶性循环的雏形。
\item 把你自己的手机充电过程当成一个黑箱：充电器标“五伏两安”，手机从零电量充到满电约一小时多一点，电池标称容量约十五瓦时。充电器一小时输出的电能约十瓦时，和电池容量对得上吗？差的那部分去了哪里？提示：注意“瓦时”是能量单位；比较时要算全程平均功率而不是最大功率；损耗主要变成充电器和手机里的热——这也是为什么快充时手机会发热。
\item 假设全世界所有铜导线里的自由电子突然约好，全部朝同一个方向以漂移速度流动，电流方向全球一致。有什么守恒律会立刻被违反？提示：想想基尔霍夫电流定律背后的那条更基本的定律；电荷不会凭空消失，它们必然在某些地方堆积——而堆积的电荷会立刻产生强大的电场把自己“推回去”。电路里我们看到的稳恒流动，是无数电荷在守恒律约束下动态平衡的结果。
\end{puzzles}
